% ============================================================================
% Fichier  : partie4/P04_C15_reseaux_protocoles.tex
% Titre    : Sécurité des Réseaux (1) : Protocoles et Attaques
% Auteur   : MINKA MI NGUIDJOÏ Thierry Emmanuel
% Version  : 1.0 -- novembre 2026
% ============================================================================
% Base : brouillon « Livre V » (modèles OSI et TCP/IP, protocoles). Sources :
% RFC 826 (ARP), RFC 4987 (inondation SYN), RFC 5452 (résistance du DNS aux
% réponses falsifiées), RFC 4033 (DNSSEC) ; Stevens, TCP/IP Illustrated.
% ============================================================================

\chapter{Sécurité des Réseaux (1) : Protocoles et Attaques}
\label{chap:reseaux1}
\chaptermark{Protocoles et attaques réseau}

\epigraph{%
    « Les protocoles d'Internet ont été conçus pour des machines\\
    qui se faisaient confiance. »%
}{Principe du cours}

\niveauChapitre{Intermédiaire}{Bases des réseaux (modèle TCP/IP,
adressage)~; Chapitre~\ref{chap:menaces}.}

\begin{boxobjectifs}
\begin{itemize}
    \item Rattacher les attaques réseau aux couches qu'elles visent.
    \item Expliquer l'usurpation ARP et l'empoisonnement DNS.
    \item Décrire l'inondation SYN et les dénis de service par
          amplification.
    \item Repérer les protocoles en clair et leurs remplaçants sûrs.
    \item Lire une capture réseau pour y trouver une attaque.
\end{itemize}
\end{boxobjectifs}

\bigskip

% ============================================================================
\section*{Question directrice : que voit celui qui écoute le réseau d'une agence~?}
% ============================================================================

Lors de l'enquête sur l'Opération GAROUA, l'équipe capture une heure de
trafic sur le réseau de l'agence de Garoua. Elle y trouve, en clair, le
mot de passe d'un serveur FTP, des requêtes d'administration en HTTP et
un balayage de ports venu d'un poste interne. Ce chapitre apprend à lire
ces traces, et à comprendre pourquoi les protocoles les laissent passer.

\medskip
\textbf{Plan du chapitre.}
\begin{enumerate}
    \item Les couches et leurs attaques (section~\ref{sec:c15-couches}).
    \item Usurper sur le réseau local~: ARP (section~\ref{sec:c15-arp}).
    \item Tromper la résolution de noms~: DNS (section~\ref{sec:c15-dns}).
    \item Saturer~: les dénis de service (section~\ref{sec:c15-dos}).
    \item Les protocoles en clair (section~\ref{sec:c15-clair}).
    \item Lire une capture (section~\ref{sec:c15-capture}).
\end{enumerate}

\separateur

% ============================================================================
\section{Les Couches et leurs Attaques}
\label{sec:c15-couches}
% ============================================================================

\begin{tableaucours}{L{2.8cm}L{3.6cm}L{5.8cm}}{%
    \tableauHeader{Couche TCP/IP} &
    \tableauHeader{Protocoles} &
    \tableauHeader{Attaques typiques}
}
Accès au réseau & Ethernet, Wi-Fi, ARP & Écoute, usurpation ARP, saturation
des tables des commutateurs \\
Internet & IP, ICMP & Usurpation d'adresse IP, fragmentation, inondation
ICMP \\
Transport & TCP, UDP & Inondation SYN, balayage de ports, détournement de
session \\
Application & DNS, HTTP, SMTP, FTP & Empoisonnement DNS, interception de
protocoles en clair, injections \\
\end{tableaucours}
\vspace{-0.8em}
\begin{center}{\small\itshape Tableau~15.1. Les attaques par couche}\end{center}

% ============================================================================
\section{Usurper sur le Réseau Local : ARP}
\label{sec:c15-arp}
% ============================================================================

Le protocole ARP associe une adresse IP à une adresse matérielle (MAC) sur
un réseau local~\cite{rfc826}. Il n'authentifie rien~: une machine accepte
une réponse ARP qu'elle n'a pas demandée.

\begin{boxdef}
\termecle{Usurpation ARP} (\textit{ARP spoofing})~: envoi de fausses
réponses ARP qui associent l'adresse IP de la passerelle à l'adresse MAC
de l'attaquant~; le trafic des victimes passe alors par lui (homme du
milieu).
\end{boxdef}

\textbf{Contre-mesures}~: inspection ARP dynamique sur les commutateurs,
segmentation en petits réseaux virtuels, et surtout chiffrement des
échanges (TLS), qui rend l'interception inutile.

% ============================================================================
\section{Tromper la Résolution de Noms : DNS}
\label{sec:c15-dns}
% ============================================================================

Le DNS traduit les noms en adresses. Un \termecle{empoisonnement de cache}
consiste à faire accepter par un résolveur une fausse réponse, qui envoie
ensuite tous ses utilisateurs vers le serveur de l'attaquant. En 2008,
D.~Kaminsky a montré qu'une attaque de ce type était bien plus facile
qu'on ne le pensait~; la réponse a été de rendre aléatoires le port
source et l'identifiant des requêtes~\cite{rfc5452}. DNSSEC va plus loin
en signant les réponses~\cite{rfc4033}.

\begin{boxalert}
\textbf{Violation typique.} Laisser un résolveur DNS interne répondre aux
requêtes venues d'Internet (résolveur ouvert).

\smallskip
\textbf{Règle~:} un résolveur ne sert que ses utilisateurs~; ouvert, il
sert aussi les attaquants, notamment pour l'amplification
(section~\ref{sec:c15-dos}).
\end{boxalert}

% ============================================================================
\section{Saturer : les Dénis de Service}
\label{sec:c15-dos}
% ============================================================================

\begin{tableaucours}{L{3.4cm}L{5.0cm}L{3.8cm}}{%
    \tableauHeader{Attaque} &
    \tableauHeader{Principe} &
    \tableauHeader{Contre-mesure}
}
Inondation SYN & Ouvrir des milliers de connexions TCP sans les achever,
pour épuiser la mémoire du serveur & Témoins SYN (\textit{SYN
cookies})~\cite{rfc4987} \\
Amplification & Envoyer de petites requêtes avec l'adresse de la victime,
à des serveurs qui répondent par de grosses réponses (DNS, NTP) &
Fermer les services ouverts, filtrer les adresses usurpées \\
Déni de service applicatif & Saturer une fonction coûteuse (recherche,
connexion) & Limitation de débit, cache \\
Déni de service distribué & Des milliers de machines (botnet) & Service
de filtrage de l'opérateur ou d'un fournisseur spécialisé \\
\end{tableaucours}
\vspace{-0.8em}
\begin{center}{\small\itshape Tableau~15.2. Les dénis de service}\end{center}

Le \termecle{facteur d'amplification} est le rapport entre la taille de
la réponse et celle de la requête. Une requête de 60~octets qui provoque
une réponse de 3\,000~octets amplifie 50~fois~: 1~Mbit/s envoyé par
l'attaquant devient 50~Mbit/s reçus par la victime.

% ============================================================================
\section{Les Protocoles en Clair}
\label{sec:c15-clair}
% ============================================================================

\begin{tableaucours}{L{2.8cm}L{4.6cm}L{4.8cm}}{%
    \tableauHeader{Protocole en clair} &
    \tableauHeader{Ce qui circule lisiblement} &
    \tableauHeader{Remplaçant}
}
Telnet & Identifiants et commandes & SSH \\
FTP & Identifiants et fichiers & SFTP, FTPS \\
HTTP & Formulaires, cookies de session & HTTPS (TLS) \\
SMTP sans TLS & Contenu des courriels & SMTP avec STARTTLS ou TLS \\
SNMP v1, v2c & Chaîne de communauté, configuration & SNMPv3 avec
authentification et chiffrement \\
POP3, IMAP sans TLS & Identifiants, courriels & POP3S, IMAPS \\
\end{tableaucours}
\vspace{-0.8em}
\begin{center}{\small\itshape Tableau~15.3. Protocoles en clair et remplaçants}\end{center}

% ============================================================================
\section{Lire une Capture}
\label{sec:c15-capture}
% ============================================================================

\begin{boxoutils}{Filtres de capture et d'affichage}
\begin{lstlisting}
# capturer le trafic d'une interface dans un fichier
tcpdump -i eth0 -w agence.pcap
# Wireshark : quelques filtres d'affichage
ftp.request.command == "USER" || ftp.request.command == "PASS"
http.request.method == "POST"
tcp.flags.syn == 1 && tcp.flags.ack == 0
arp.duplicate-address-detected
\end{lstlisting}
Le troisième filtre isole les débuts de connexion~: un grand nombre de
SYN vers des ports différents d'une même machine signe un balayage.
\end{boxoutils}

\begin{boxscenario}{Opération GAROUA. Une heure de trafic à Garoua}
\textbf{Constatations} (capture du 10~mars, 9~h à 10~h)~:
\begin{itemize}
    \item 4 connexions FTP vers le serveur de dépôt des bordereaux~;
          identifiant et mot de passe lisibles en clair~;
    \item des requêtes HTTP vers l'interface d'administration d'une
          imprimante réseau, avec le mot de passe par défaut~;
    \item depuis le poste PC-GAR-07, des paquets SYN vers 1\,000~ports de
          chacun des 12~serveurs, en moins de deux minutes.
\end{itemize}

\textbf{Analyse} (confiance moyenne). Le balayage depuis PC-GAR-07 laisse
penser que ce poste est encore compromis, ou utilisé par un outil non
autorisé~: il faut l'isoler et l'examiner avant de conclure. Les mots de
passe FTP et de l'imprimante doivent être changés et les protocoles
remplacés.
\end{boxscenario}

\begin{boxmetier}{Analyste SOC}
\textbf{Reconnaître les motifs.} L'analyste reconnaît un balayage (une
source, de nombreux ports), une inondation SYN (de nombreuses sources, des
connexions jamais achevées) et une usurpation ARP (deux adresses MAC pour
la même adresse IP).
\end{boxmetier}

\begin{boxmetier}{Administrateur systèmes et réseaux}
\textbf{Éteindre le clair.} L'administrateur inventorie les services en
clair, les remplace, et bloque leurs ports aux frontières des zones.
\end{boxmetier}

\separateur

\begin{boxresume}
\textbf{Ce qu'il faut retenir de ce chapitre~:}
\begin{itemize}
    \item \textbf{Chaque couche} a ses attaques~; les protocoles anciens
          n'authentifient rien.
    \item \textbf{L'usurpation ARP} place l'attaquant au milieu~; le
          chiffrement la rend inutile.
    \item \textbf{L'empoisonnement DNS} se combat par l'aléa des requêtes
          et par DNSSEC~; un résolveur ne doit pas être ouvert.
    \item \textbf{Les dénis de service} épuisent ou amplifient~; les
          témoins SYN et le filtrage les limitent.
    \item \textbf{Les protocoles en clair} ont tous un remplaçant chiffré.
    \item \textbf{Une capture} se lit avec des filtres~; une constatation
          réseau appelle souvent une vérification sur la machine.
\end{itemize}
\end{boxresume}

\begin{boxinfo}
\textbf{Points de contrôle de l'élève.} Je sais~: rattacher une attaque à
sa couche~; expliquer l'usurpation ARP et ses contre-mesures~; calculer un
facteur d'amplification~; citer le remplaçant de chaque protocole en
clair~; écrire un filtre Wireshark pour un balayage.
\end{boxinfo}

\bigskip

\begin{boxexo}[title={\ding{45}\quad Exercice~15.1. Quelle couche~? \niveaubase}]
Rattachez à une couche~: inondation SYN~; empoisonnement DNS~; usurpation
ARP~; écoute du Wi-Fi~; injection SQL.
\end{boxexo}

\begin{boxexo}[title={\ding{45}\quad Exercice~15.2. Remplacer \niveaubase}]
Donnez le remplaçant sûr de Telnet, FTP, HTTP et SNMPv2c, et le port
habituel de chacun.
\end{boxexo}

\begin{boxexo}[title={\ding{45}\quad Exercice~15.3. Amplification \niveaubase}]
Une requête de 40~octets provoque une réponse de 4\,000~octets. Quel est
le facteur d'amplification~? Quel débit reçoit la victime si l'attaquant
envoie 2~Mbit/s~?
\end{boxexo}

\begin{boxexo}[title={\ding{45}\quad Exercice~15.4. ARP \niveaumoyen}]
Dans une table ARP, deux adresses IP différentes ont la même adresse MAC.
Qu'en déduisez-vous~? Quelles vérifications faites-vous~?
\end{boxexo}

\begin{boxexo}[title={\ding{45}\quad Exercice~15.5. SYN cookies \niveaumoyen}]
Expliquez pourquoi les témoins SYN protègent contre l'inondation SYN sans
conserver d'état pour les connexions incomplètes.
\end{boxexo}

\begin{boxexo}[title={\ding{45}\quad Exercice~15.6. Filtres \niveaumoyen}]
Écrivez les filtres Wireshark qui isolent~: les requêtes DNS vers un
serveur extérieur à l'entreprise~; les connexions Telnet~; les réponses
ICMP « port inaccessible ».
\end{boxexo}

\begin{boxexo}[title={\ding{45}\quad Exercice~15.7. Le balayage \niveaumoyen}]
Reprenez le scénario de Garoua. Rédigez séparément la constatation et
l'analyse sur PC-GAR-07, et les trois actions à mener.
\end{boxexo}

\begin{boxexo}[title={\ding{45}\quad Exercice~15.8. Le résolveur ouvert \niveaumoyen}]
Le résolveur DNS de l'ANPP répond aux requêtes venues d'Internet. Quels
risques pour l'ANPP et pour les autres~? Quelle correction~?
\end{boxexo}

\begin{boxexo}[title={\ding{45}\quad Exercice~15.9. Un plan de chiffrement \niveauexpert}]
Proposez à LiZbiZ un plan pour supprimer en six mois tous les protocoles
en clair~: inventaire, priorités, remplacement, contrôle.
\end{boxexo}

\begin{boxexo}[title={\ding{45}\quad Exercice~15.10. Le portail sous DDoS \niveauexpert}]
Le portail de suivi de LiZbiZ subit un déni de service distribué la
veille des fêtes. Rédigez la procédure de réponse et les mesures
préventives à négocier avec l'opérateur.
\end{boxexo}

\bigskip

\begin{boxinfo}
\textbf{Pour aller plus loin.}
\begin{itemize}
    \item RFC~826 (ARP), RFC~4987 (inondation SYN), RFC~5452 et RFC~4033
          (DNS)~\cite{rfc826,rfc4987,rfc5452,rfc4033}.
    \item W.~R.~Stevens, \emph{TCP/IP Illustrated}~\cite{stevens1994}.
\end{itemize}
\end{boxinfo}

\begin{boxinfo}
\textbf{Transition.} Connaître les attaques ne suffit pas~: il faut les
arrêter et les voir. Le Chapitre~\ref{chap:reseaux2} présente les
équipements et les protocoles qui protègent le réseau.
\end{boxinfo}
